\begin{frame}[fragile]{Configuraciones iniciales 1/2}

    \begin{block}{Tu identidad}
        Es importante establecer nuestro \textbf{nombre y email} en nuestro repositorio, ya que estos van a ir asociados con los cambios que hagamos:

        \vspace{0.5em}

        \texttt{git config --global user.name "Guybrush Threepwood"}

        \texttt{git config --global user.email guybrush@example.com}
    \end{block}

    \begin{block}{Config global vs local}
    También se puede omitir el flag \texttt{--global} para que las credenciales apliquen solo al repositorio en el que estamos trabajando.
    \end{block}

\end{frame}
\begin{frame}[fragile]{Configuraciones iniciales 2/2}

	\begin{block}{Autenticación con el servidor}
        Para poder acceder al repositorio remoto, el servidor necesita saber si tenemos permisos para accederlo. Para eso, cada vez que ejecutemos \texttt{git pull} o \texttt{git push} vamos a tener que completar nuestro usuario y contraseña.

        %De esta manera, cuando tengamos un repositorio en Gitlab vamos a poder conectarlo con nuestro local.
	\end{block}

	\begin{block}{SSH}
        Para no tener que estar escribiendo constantemente nuestra contraseña, podemos configurar otras formas de autenticación como SSH. En esta versión corta del taller no llegaremos a verlo.

        %De esta manera, cuando tengamos un repositorio en Gitlab vamos a poder conectarlo con nuestro local.
	\end{block}
\end{frame}
